\begin{textsample}{NEG Dim 5 – human – Score -9.00 – rj/rj\_ef\_intro\_7\_p1.txt}  \label{ex:f5_neg_018}
As dez competências gerais da BNCC, definidas a partir dos direitos éticos, estéticos e políticos garantidos pelas Diretrizes \textbf{Curriculares} Nacionais, apresentam -se como uma indicação clara do que os alunos precisam “ saber ” e “ saber fazer ”.
Essas competências são desdobradas em competências específicas de cada área do \textbf{conhecimento} do Ensino Fundamental que, por sua vez, são formadas por um ou mais componentes \textbf{curriculares}, sendo que cada um possui competências específicas de componentes.
Ao \textbf{longo} dos \textbf{anos}, cada componente deve apresentar um conjunto de \textbf{habilidades} mais complexo, para garantir o \textbf{desenvolvimento} dos alunos.
Essas \textbf{habilidades}, por sua vez, são \textbf{relacionadas} a diferentes \textbf{objetos} de \textbf{conhecimento}, organizadas em unidades temáticas.

% matched lemmas: ano, conhecimento, curricular, desenvolvimento, habilidade, longo, objeto, relacionar
\end{textsample}
